\documentclass[a4paper,
fontsize=11pt,
%headings=small,
oneside,
numbers=noperiodatend,
parskip=half-,
bibliography=totoc,
final
]{scrartcl}

\usepackage[babel]{csquotes}
\usepackage{synttree}
\usepackage{graphicx}
\setkeys{Gin}{width=.4\textwidth} %default pics size

\graphicspath{{./plots/}}
\usepackage[ngerman]{babel}
\usepackage[T1]{fontenc}
%\usepackage{amsmath}
\usepackage[utf8x]{inputenc}
\usepackage [hyphens]{url}
\usepackage{booktabs} 
\usepackage[left=2.4cm,right=2.4cm,top=2.3cm,bottom=2cm,includeheadfoot]{geometry}
\usepackage[labelformat=empty]{caption} % option 'labelformat=empty]' to surpress adding "Abbildung 1:" or "Figure 1" before each caption / use parameter '\captionsetup{labelformat=empty}' instead to change this for just one caption
\usepackage{eurosym}
\usepackage{multirow}
\usepackage[ngerman]{varioref}
\setcapindent{1em}
\renewcommand{\labelitemi}{--}
\usepackage{paralist}
\usepackage{pdfpages}
\usepackage{lscape}
\usepackage{float}
\usepackage{acronym}
\usepackage{eurosym}
\usepackage{longtable,lscape}
\usepackage{mathpazo}
\usepackage[normalem]{ulem} %emphasize weiterhin kursiv
\usepackage[flushmargin,ragged]{footmisc} % left align footnote
\usepackage{ccicons} 
\setcapindent{0pt} % no indentation in captions
\usepackage{xurl} % Breaks URLs

%%%% fancy LIBREAS URL color 
\usepackage{xcolor}
\definecolor{libreas}{RGB}{112,0,0}

\usepackage{listings}

\urlstyle{same}  % don't use monospace font for urls

\usepackage[fleqn]{amsmath}

%adjust fontsize for part

\usepackage{sectsty}
\partfont{\large}

%Das BibTeX-Zeichen mit \BibTeX setzen:
\def\symbol#1{\char #1\relax}
\def\bsl{{\tt\symbol{'134}}}
\def\BibTeX{{\rm B\kern-.05em{\sc i\kern-.025em b}\kern-.08em
    T\kern-.1667em\lower.7ex\hbox{E}\kern-.125emX}}

\usepackage{fancyhdr}
\fancyhf{}
\pagestyle{fancyplain}
\fancyhead[R]{\thepage}

% make sure bookmarks are created eventough sections are not numbered!
% uncommend if sections are numbered (bookmarks created by default)
\makeatletter
\renewcommand\@seccntformat[1]{}
\makeatother

% typo setup
\clubpenalty = 10000
\widowpenalty = 10000
\displaywidowpenalty = 10000

\usepackage{hyperxmp}
\usepackage[colorlinks, linkcolor=black,citecolor=black, urlcolor=libreas,
breaklinks= true,bookmarks=true,bookmarksopen=true]{hyperref}
\usepackage{breakurl}

%meta
%meta

\fancyhead[L]{C. Lenz, M. Miller\\ %author
LIBREAS. Library Ideas, 48 (2026). % journal, issue, volume.
\href{https://doi.org/10.18452/37903}{\color{black}https://doi.org/10.18452/37903}
{}} % doi
\fancyhead[R]{\thepage} %page number
\fancyfoot[L] {\ccLogo \ccAttribution\ \href{https://creativecommons.org/licenses/by/4.0/}{\color{black}Creative Commons BY 4.0}}  %licence
\fancyfoot[R] {ISSN: 1860-7950}

\title{\LARGE{Interviews als Basis des Wissenstransfers }}% title
\author{Charlotte Lenz, Matthias Miller} % author

\setcounter{page}{1}

\hypersetup{%
      pdftitle={Interviews als Basis des Wissenstransfers },
      pdfauthor={Charlotte Lenz, Matthias Miller},
      pdfcopyright={CC BY 4.0 International},
      pdfsubject={LIBREAS. Library Ideas, 48 (2026)},
      pdfkeywords={Interviews, Übergabe Bibliothek, Bibliotheksmanagement, Wissenstransfer, Nachfolge, library management, knowledge transfer},
      pdflicenseurl={https://creativecommons.org/licenses/by/4.0/},
      pdfcontacturl={http://libreas.eu},
      pdfurl={https://doi.org/10.18452/37903},
      pdfdoi={10.18452/37903},
      pdflang={de},
      pdfmetalang={de}
     }



\date{}
\begin{document}

\maketitle
\thispagestyle{fancyplain} 

%abstracts
\begin{abstract}
\noindent
\textbf{Kurzfassung}: Neben Anwendung als wissenschaftliche
Forschungsmethode können Expert\*inneninterviews in Bibliotheken als
Möglichkeit herangezogen werden, um vorhandenes Fachwissen zu
dokumentieren, damit ein aktueller reibungsloser Arbeitsablauf im
Betrieb gewährleistet werden kann. Zudem gestattet solch eine
Dokumentation neuen Mitarbeiter\*innen, sich möglichst schnell im neuen
Wirkungsbereich zurechtzufinden. Eine weitere Anwendung finden
Interviews in Bibliotheken für den Erfahrungsaustausch unter
Fachkolleg\*innen. Die Autor\*innen testen als ehemalige direkte
Kolleg*innen letztere Variante zum Thema Neukonzeption von Bibliotheken
im Zuge eines Arbeitsstellenwechsels. Der besseren Lesbarkeit willen
erfolgen alle weiteren personenbezogenen Bezeichnungen in diesem Aufsatz
in der männlichen Form, welche aber keinerlei wertende Bedeutung hat und
alle Geschlechter gleichermaßen repräsentiert.

\begin{center}\rule{0.5\linewidth}{0.5pt}\end{center}

\noindent\textbf{Abstract}: Expert interviews in libraries can be used as a
scientific research method, but also for documenting existing expertise
for ensuring smooth and efficient workflows. Such documentation also
allows new employees to quickly familiarize themselves with their new
work environment. Another application of interviews in libraries is the
exchange of experiences among colleagues. As former direct colleagues,
the authors are testing this latter approach in the context of library
redesign following off a job change.
\end{abstract}

%body
\section{Einleitung}\label{einleitung}

Experteninterviews als wissenschaftlich anerkannte Forschungsmethode
sind mittlerweile auch in der Bibliothekswissenschaft eine feste Größe.
Vielen Forschungsarbeiten liegen daher strukturierte oder
halbstrukturierte, meist anhand eines Fragenkatalogs geführte Gespräche
zugrunde. Ziel ist zum einen, das zum Teil hochspezialisierte Wissen von
Experten eines Fachgebiets zu sammeln und zu strukturieren (Vgl.
Pfadenhauer, 2002, S. 113), zum anderen aber auch, aus den
Erkenntnissen, die aus diesen Gesprächen gewonnen werden, logische
Schlüsse für eine mögliche individuelle Forschungsfrage zu ziehen und
damit Leitlinien für die Praxis abzuleiten. Voraussetzung dafür ist,
dass es sich bei diesem Expertenwissen nicht nur um eine übergreifende
Kenntnis zu einem einzelnen Bereich handelt, sondern wenn möglich um ein
besonderes Spezialwissen (Vgl. Pfadenhauer, 2002, S. 115), dessen
Grundlage die Kompetenz des befragten Experten ist. Dies bedeutet, dass
er auf der einen Seite gesondertes Wissen zu der Materie haben sollte,
mit der sich der Interviewführende beschäftigt, auf der anderen Seite
aber in seiner tatsächlichen Praxis auch über eine gewisse Zuständigkeit
verfügen sollte, um Lösungen für mit der Materie einhergehende Probleme
aufzuzeigen beziehungsweise im besten Fall sogar bereitzustellen. (Vgl.
Pfadenhauer, 2002, S. 116)

Eine besondere Form des Experteninterviews stellt ein Austausch dar, bei
dem sowohl der Fragende als auch der Befragte vom Fach, ergo beide
Experten, sind. Diese Gespräche dienen der gegenseitigen Information und
dem Ausbau des jeweiligen Fachwissens zu einem bestimmten Thema, was
aber in keinem Fall als eine Art Belehrung verstanden wird, sondern als
eine objektive Darstellung, wie zum Beispiel bei einer bestimmten
Angelegenheit vorgegangen werden kann. Nicht selten werden die
aufgezeigten Vorgehensweisen sachlich diskutiert, eben weil es sich um
ein Gespräch auf Augenhöhe unter gleichberechtigten Partnern handelt.
(Vgl. Pfadenhauer, 2002, S. 119)

Für Fachpersonal in Bibliotheken eignen sich gezielte Befragungen ebenso
wie geplante Dialoge im Sinne des zuletzt erläuterten Austauschs
zwischen gleichwertigen Experten für das Wissensmanagement in der
eigenen Einrichtung und dem Teilen von Wissen bezüglich vergleichbarer
Situationen und Sachverhalte. Bei Ersterem kann zum Beispiel
detailliertes Wissen, welches sich ein Mitarbeiter individuell über
Jahre selbst bezüglich seines Tätigkeitsbereiches angeeignet hat,
gesammelt, dokumentiert und für eine nachfolgende Anwendung gespeichert
werden, oder dieses Spezialwissen wird direkt über ein Gespräch an eine
gleichrangige Person weitergegeben. Bei Letzterem wird das Interview
geführt, um aus den Erfahrungen, die eine vergleichbare Person bereits
gemacht hat, zu profitieren und für die eigene Situation anwendbare
Empfehlungen herzuleiten. Diese Art des Interviews wurde von den Autoren
des Artikels getestet. Die folgenden Ausführungen stellen den Kontext
dar, in dem das Gespräch geplant wurde, dokumentieren dessen
Vorbereitung, Ablauf und Ergebnis und beschreiben dann eingehend den
Erkenntniswert des Gesprächs aus Sicht des Befragten und des
Gesprächsleiters.

\section{Ausgangssituation}\label{ausgangssituation}

Mehr als 17 Jahre waren die Autoren direkte Kollegen in einer Berliner
Museumsbibliothek. Während der Befragte von Anfang an die Stellung als
Bibliotheksleiter innehatte, war die Interviewführende dort zunächst als
Fachangestellte und später als wissenschaftliche Bibliothekarin tätig.
Zuletzt wurde der Befragte durch die Interviewführende in Abwesenheit
vertreten. Im letzten Jahr verließ Letztere auf eigenen Wunsch die
Museumsbibliothek und übernahm die Verantwortung für eine kleine
Gedenkstättenbibliothek in Form einer One-Person-Library. Da sich diese
Bibliothek heute konzeptionell und technisch in einem ähnlichen Zustand
befindet, wie die Museumsbibliothek zu dem Zeitpunkt, als der Befragte
dort Leiter wurde, wurde die Gelegenheit für ein formelles Gespräch über
eine Neukonzeptionierung der Gedenkstättenbibliothek geschaffen.

\section{Vorbereitung und
Durchführung}\label{vorbereitung-und-durchfuxfchrung}

Nach dem Vereinbaren eines Termins, der auf eineinhalb Stunden
festgesetzt wurde, arbeitete die Interviewführende einen Katalog mit 10
Leitfragen aus, der dem Befragten vorab zur Vorbereitung zur Verfügung
gestellt wurde. Die Fragen gaben dabei den Rahmen des Gesprächs vor und
sollten sicherstellen, dass kein wichtiger Aspekt der behandelten
Thematik vergessen wird, ließen aber dennoch Raum für einen offenen,
weiterführenden Dialog. Das Gespräch selbst wurde als Onlinetermin über
die Kommunikationsplattform Microsoft Teams organisiert und während der
Durchführung von diesem System aufgenommen und im Nachhinein
transkribiert. Obwohl die Technik in solch einem Zusammenhang von beiden
Seiten erstmalig erprobt wurde, lief der Termin diesbezüglich
reibungslos ab.

Nach einer Begrüßung erklärte die Interviewführende den Ablauf und das
Ziel des Gesprächs, holte sich die Zustimmung für die Aufzeichnung der
Unterhaltung ein und stellte dar, warum sie in dem Befragten einen
passenden Interviewpartner für das Thema Neukonzeption einer Bibliothek
sieht. Danach bat sie ihn, obwohl sie ihn persönlich kannte, sich selbst
noch einmal vorzustellen. Dies entspannte die Situation erheblich, da es
in der Regel nichts Einfacheres gibt, als über sich selbst und das, was
man den ganzen Tag tut, zu sprechen.

Im Anschluss erfolgte die Befragung anhand der zuvor vorbereiteten
Fragen. Diese wurden vom Befragten der Reihe nach systematisch
beantwortet, wobei persönliche Erfahrungen und Sichtweisen ebenso wie
theoretisches Basiswissen mit in die Antworten einflossen. Die
Interviewführende bemühte sich, die Informationen parallel zu
reflektieren und auf ihre eigene Situation anzuwenden. Da gleich zu
Beginn feststand, dass das Gespräch neben dem Erfahrungsaustausch auch
als Basis für einen Fachartikel verwendet werden soll, wurde es mit
einem kurzen Austausch über das weitere Vorgehen und das Verwenden
personenbezogener beziehungsweise hausinterner Daten abgeschlossen.

\section{Allgemein-organisatorischer Erkenntniswert aus Sicht der
Interviewführenden}\label{allgemein-organisatorischer-erkenntniswert-aus-sicht-der-interviewfuxfchrenden}

Insgesamt gesehen ist das Gespräch, das zudem einen Erstversuch
darstellte, aus Sicht der Interviewführenden mehr als positiv verlaufen.
Die inhaltlichen und technischen Vorbereitungen zeigten sich als
ausreichend. Die Technik überzeugte vor allem in der anschließenden
unkomplizierten Transkription des Interviews, was eine spätere
schriftliche Auswertung enorm vereinfachte. Da für die Interviewführende
das Gespräch gleichzeitig einen Test für zukünftige Experteninterviews
im Rahmen ihrer Promotion darstellte, war gerade das Ausprobieren der
Technik ein wichtiger Faktor. Zudem konnte sie darüber hinaus wertvolle
Erkenntnisse zum Ablauf des Gesprächs für sich ableiten, so zum
Beispiel, dass das Zurverfügungstellen des Fragenkatalogs an den
Befragten im Vorfeld zu einem flüssigen Gesprächsverlauf führt, dass das
anvisierte Gesprächszeitfenster für bis zu 10 Fragen ausreichend ist
sowie dass die Eröffnungssequenz nicht nur eine angenehme Atmosphäre
schafft, sondern auch bis zu 10 Minuten der Gesamtzeit in Anspruch
nehmen kann. Gerade bei einem Gespräch mit jemand Bekanntem, der zu
großen Teilen das Gleiche erlebt hat, kommt der Wert der anderen
Perspektive auf ein und dasselbe Erlebnis mit hinzu. Ob jedes Interview
in dieser Form so ablaufen kann, ist allerdings fraglich, da man als
Interviewführender nicht mit jedem Befragten eine gemeinsame
Arbeitsvergangenheit teilt. Zudem sollte eine simultane Reflexion des
Gesagten auf die eigene Situation in einem Expertengespräch, welches als
Forschungsmethode für eine wissenschaftliche Untersuchung durchgeführt
wird, zumindest laut unterlassen werden, weil es bei diesem darum geht,
zunächst objektive Daten zu erheben und den Befragten nicht mit eigenen
Gedanken und Schlussfolgerungen zu einem bestimmten Meinungsbild hin zu
beeinflussen.

\section{Inhaltlicher Erkenntniswert aus Sicht der
Interviewführenden}\label{inhaltlicher-erkenntniswert-aus-sicht-der-interviewfuxfchrenden}

Der größte Mehrwert eines Interviews liegt zweifellos in der
inhaltlichen Erkenntnisgewinnung. Bezogen auf den thematischen Rahmen --
die Neukonzeption einer Bibliothek im Rahmen eines Arbeitswechsels --
konnte die Interviewführende folgende 10 wertvolle Informationen aus dem
fachlichen Austausch ziehen:

\begin{enumerate}
\def\labelenumi{\arabic{enumi}.}
\item
  Auch mit vorliegender, zum Teil langjähriger Berufserfahrung,
  gestaltet sich ein Neuanfang in einer anderen Einrichtung nicht
  einfach. Ist die neue Position in einer Management- oder
  Leitungsfunktion angesiedelt, kann dieser sich sogar als sehr
  anstrengend erweisen.
\item
  Die Theorie bereitet einen nicht vollumfänglich auf die Praxis vor.
  Die Ausbildung im Bibliotheksbereich kann die Basis für eine
  zukünftige Anstellung bilden, auf der allerdings mit eigenen
  Erfahrungen aufgebaut werden muss. Viele Dinge werden in der
  Ausbildung nämlich entweder nicht oder nur unzureichend vermittelt,
  beziehungsweise es wird der Idealfall gelehrt, der in der Praxis
  aufgrund von statischen Verwaltungsabläufen und individuellen
  Einflüssen der Akteure aber nicht standhält.
\item
  In der ersten Zeit auf der neuen Position ist es enorm wichtig, sich
  mit organisatorischen Dingen und Verwaltungsabläufen vertraut zu
  machen. Diese sind in Einrichtungen des öffentlichen Dienstes oft
  ähnlich, allerdings hat jeder Betrieb auch eigene gesonderte Prozesse.
  Diese zu kennen, hilft, Ärgernisse für andere Bereiche zu vermeiden,
  gerade an jenen wichtigen Stellen, von denen man sich in Zukunft bei
  der Umsetzung von Projekten et cetera Hilfe erhofft.
\item
  Die Neukonzeption einer Bibliothek geht immer mit dem Setzen von
  konkreten Zielen einher. Diese können sich in völlig unterschiedlichen
  Bereichen zeigen, betreffen allerdings bei kleineren Bibliotheken in
  der Regel die technische Weiterentwicklung, welche meist auch die
  überregionale Sichtbarkeit der Bestände impliziert, sowie eine
  institutionelle und übergreifende Vernetzung, weil man als kleine
  Einrichtung auf die Unterstützung anderer (möglicherweise größerer
  Player) angewiesen ist.
\item
  Eine Einarbeitung und Führung durch den übergeordneten Leiter ist bei
  Leitungspositionen eher nicht zu erwarten. Gewöhnlich muss man relativ
  schnell eigenverantwortlich klarkommen und sich anderweitig Hilfe
  suchen.
\item
  Aufgrund dessen sollte man für die Amtshilfe von Kollegen der gleichen
  Hierarchieebene stets offen und dankbar sein. Unterstützung kann sogar
  von Kollegen erfolgen, die einem unterstellt sind, weil sie angesichts
  einer möglichen langjährigen Institutionszugehörigkeit über spezielles
  Wissen verfügen, das man selbst noch nicht hat. Zudem fördert dies den
  Teamgeist, der für zukünftige Zielsetzungen und Konzepte essentiell
  ist.
\item
  Die Neukonzeption einer Bibliothek ist immer nur ein erster Schritt,
  denn die Bibliothekswelt ist im stetigen Wandeln. Dies führt dazu,
  dass man über die Zeit mit Themen, Aufgaben und Problemen konfrontiert
  wird, die man zu Beginn noch nicht im Blick hatte oder auch noch nicht
  haben konnte. Somit kommt, auch wenn man schon viel Erfahrung in
  seinem Beruf hat, immer wieder Neues hinzu, womit es sich
  auseinanderzusetzen gilt.
\item
  Mit Neukonzeptionen gehen auch immer Projekte einher, die
  höhergestellten Entscheidungsträgern, welche meist nicht vom Fach
  sind, erklärt werden müssen. Eine gute Argumentationskultur sollte man
  sich daher als Kernkompetenz aneignen.
\item
  Nichts geht über eine gute institutionsinterne und
  institutionsübergreifende Vernetzung. Gerade kleinere Bibliotheken
  sind auf Kooperationen mit anderen angewiesen, um für Innovationen die
  fachliche und technische Unterstützung von anderen Abteilungen der
  eigenen Einrichtung sowie die von externen Partnern zu gewinnen. Diese
  können sich positiv auf personelle und finanzielle Ressourcen bei der
  Umsetzung auswirken.
\item
  Freue Dich über Teilerfolge und übe Dich in Geduld, denn keine
  Neukonzeption geht ohne Widerstände einher. Tatsächlich findet ein
  Wandel nie ohne Hürden statt. Im Gegenteil: Für viele Dinge braucht
  man Ausdauer, da gerade Verwaltungsabläufe teilweise nur langsam
  vorangehen und Entscheidungsträger bei der Abarbeitung von an sie
  gestellten Anträgen priorisieren. Oft rutschen dabei
  Bibliotheksangelegenheiten weit nach unten. Beharrlichkeit ist an
  dieser Stelle gefragt, ebenso die Erkenntnis, dass manche Ziele
  möglicherweise nie oder nicht so umgesetzt werden können, wie man es
  sich erhofft hat.
\end{enumerate}

Abschließend wurde beiderseitig festgestellt, dass der Wechsel auf eine
neue Arbeitsstelle, welche die Neuaufstellung einer Bibliothek und ihrer
Services beinhaltet, in der Mitte des Arbeitslebens als guter Zeitpunkt
angesehen werden kann. Man hat noch entsprechend Zeit, um neue Dinge
voranzubringen, blickt aber dennoch schon auf eine gewisse
Berufserfahrung zurück, welche einem in vielen Bereichen eine gewisse
Sicherheit für den neuen Job und die damit verbundenen Herausforderungen
gibt.

\section{Der Erkenntniswert aus Sicht des
Befragten}\label{der-erkenntniswert-aus-sicht-des-befragten}

Aus Sicht des Befragten bei einem Experteninterview ist zunächst
festzustellen, dass es eine spannende Erkenntnis ist, überhaupt als
Experte wahrgenommen und für das Interview ausgewählt worden zu sein.
Schon beim ersten Gespräch mit der Interviewerin war die ernsthafte
Vorbereitung und das klare Interesse, das Thema aus der Praxis heraus zu
beleuchten, zu beobachten. Das Gespräch bot dem Befragten die
Gelegenheit, eigene Erfahrungen und Beobachtungen zu reflektieren und
gleichzeitig Impulse zu geben, die über seine bisherigen täglichen
Routinen hinausgehen.

Besonders positiv hat der Befragte die offene und respektvolle
Atmosphäre des Interviews wahrgenommen. Durch die vorbereitenden
Informationen hatte er jederzeit das Gefühl, dass seine Perspektive
geschätzt wurde und seine Beispiele tatsächlich für die Analyse des
Themas relevant waren. Es war sehr motivierend, rückblickend detailliert
über die Projekte und Strategien zu sprechen, die in seiner Einrichtung
entwickelt wurden, um die Museumsbibliothek von einem Rohdiamanten in
einen funkelnd geschliffenen Edelstein zu verwandeln. Dass auf diese
Weise dem Wissenstransfer nach dem Stellenwechsel der Interviewführenden
Vorschub geleistet wurde, war eines der Hauptanliegen bei dem Interview.

Eine weitere positive Erfahrung war die Gelegenheit, über die weichen
Faktoren zu sprechen, die oft unterschätzt werden. Der Befragte konnte
erklären, wie Vertrauen, Motivation und persönliche Beziehungen den
Erfolg von Projekten und Umgestaltungsprozessen entscheidend
beeinflussen. Dass die Interviewerin gezielt nach diesen Aspekten
fragte, gab ihm die Chance, auch subtile, nicht formal dokumentierte
Prozesse zu beleuchten, die in der Praxis häufig den größten Unterschied
machen. Dabei wurde dem Befragten deutlich, wie viel unbewusstes Wissen
in einer Einrichtung existiert und wie wichtig es ist, dieses aktiv zu
sichern, wenn Mitarbeitende die Einrichtung verlassen.

Das Interview bot dem Befragten zudem einen wertvollen Reflexionsraum.
Durch die gezielten Nachfragen konnten Situationen analysiert werden,
die sonst vielleicht nur intuitiv gehandhabt worden wären. Besonders
spannend war die Diskussion über die langfristigen Effekte von
Umgestaltungsprozessen und technischer Erneuerung in Bibliotheken. Beim
Thema Wissenstransfer zeigte sich, dass sich nur gut dokumentiertes und
kommuniziertes Wissen transferieren und auf eine neue Stelle übertragen
lässt, und sich dieses auf die Innovationsfähigkeit, die Teamdynamik und
die Stabilität von Prozessen in der neuen Einrichtung überführen lassen.
Es war befriedigend zu sehen, dass die Erfahrungen und Empfehlungen des
Befragten potentiell dazu beitragen können, dass andere Einrichtungen
von ähnlichen Strategien profitieren.

Insgesamt war das Experteninterview eine durchweg positive Erfahrung. Es
bot nicht nur die Gelegenheit, Wissen weiterzugeben, sondern auch selbst
neues Verständnis für die Bedeutung und die Mechanismen des
Wissenstransfers zu entwickeln. Der Befragte ging aus dem Gespräch mit
dem Gefühl, dass durch dieses Interview sowohl seine eigene Praxis
reflektiert und gestärkt wurde, als auch ein Beitrag zu einem besseren
Verständnis des Themas geleistet werden konnte. Es hat gezeigt, wie
wichtig es ist, Wissen bewusst zu teilen und Strukturen zu schaffen, die
den Transfer auch über institutionelle Grenzen hinweg nachhaltig
unterstützen.

\section{Zusammenfassung}\label{zusammenfassung}

Ein Experteninterview zum Thema Wissenstransfer beim Stellenwechsel kann
für eine Bibliothek auf mehreren Ebenen einen wichtigen Beitrag leisten,
um den Erhalt von Wissen zu gewährleisten. Zunächst dient es der
Identifikation kritischen Wissens. Durch gezielte Fragen kann das
Interview aufdecken, welches Wissen für den reibungslosen Betrieb und
die Kontinuität besonders entscheidend ist. In einer Bibliothek kann das
etwa spezielles Fachwissen über Katalogsysteme, Archivierungsprozesse,
Sondersammlungen oder interne Geschäftsgänge sein. Experteninterviews
helfen, diese \enquote{unsichtbaren Wissensschätze} zu erkennen, die
sonst oft nur im Kopf einzelner Mitarbeitender existieren. Gleichzeitig
kann ein Experteninterview der Dokumentation informeller Wissensprozesse
Vorschub leisten. In Bibliotheken gibt es viele Abläufe, die nicht
formal dokumentiert sind -- etwa wie Mitarbeitende spezielle
Rechercheaufträge bearbeiten, welche Routinen für die Bestandskontrolle
sinnvoll sind, wie technologische Hilfsmittel optimal genutzt werden
oder wie in bestimmten Nutzungssituationen verfahren wird. Das Interview
bietet die Gelegenheit, diese informellen Prozesse systematisch zu
erfassen, so dass sie auch nach einem Stellenwechsel oder dem
Ausscheiden von Mitarbeitenden erhalten bleiben, oder scheidenden
Mitarbeitenden mitgegeben werden können. Ein weiterer wichtiger Punkt
ist die Reflexion über bestehende Transfer- und Veränderungsprozesse.
Experteninterviews regen dazu an, Veränderungsprozesse kritisch zu
hinterfragen. Der Experte kann selbstreflektierend bewerten, welche
Maßnahmen in seiner Bibliothek bereits effektiv sind, welche Maßnahmen
gut oder schlecht umgesetzt wurden und wo noch Lücken bestehen. Dies
ermöglicht der Bibliothek, gezielt Maßnahmen zu verbessern oder neue
Mechanismen einzuführen, wie etwa strukturierte Mentoring-Programme oder
digitale Wissensplattformen. Die Förderung einer Wissenskultur kann
durch diese Selbstreflexion gestärkt werden. Durch das Interview wird
das Bewusstsein für die Bedeutung des Wissenstransfers angeregt. Der
Befragte erkennt, dass sein Fachwissen wertvoll ist und systematisch
gesichert werden sollte. Dies kann auch zu einer stärkeren Bereitschaft
führen, Wissen aktiv zu teilen, was wiederum die langfristige Stabilität
der Bibliotheksprozesse unterstützt. Interviews ermöglichen es darüber
hinaus, erfolgreiche Strategien aus der Praxis zu dokumentieren und zu
verbreiten. So können andere Mitarbeitende von bewährten Methoden
lernen, ohne dass Wissen verloren geht, wenn jemand die Einrichtung
verlässt oder wenn jemand das Wissen mitnehmen möchte. Das kann auch
dazu beitragen, dass ähnliche Herausforderungen bei zukünftigem
Personalwechsel schneller und effektiver gelöst werden. Zuletzt können
Experteninterviews auch dazu verwendet werden, den Wissenstransfer
institutsübergreifend abzusichern, wenn aus dem wertvollen Wissen und
den Erfahrungen der Experten, neues Wissen und neue Erfahrungen in einer
vergleichbaren Einrichtung entstehen.

Indem Experteninterviews Wissen systematisch erfassen, werden sie zu
einem Baustein für die nachhaltige Entwicklung der Bibliothek. Sie
gewinnt Stabilität, Flexibilität und Resilienz, da personengebundenes
Wissen nicht mehr ausschließlich an Einzelne gebunden ist, sondern in
der Institution selbst verankert wird. Ein Experteninterview ist damit
nicht nur ein Mittel zur Informationssammlung, sondern ein strategisches
Instrument für den Wissenserhalt. Es hilft, kritisches und implizites
Wissen zu identifizieren, zu dokumentieren und in die Abläufe der
täglichen Arbeit zu überführen. Für Bibliotheken bedeutet dies, dass
Abläufe, Spezialwissen und Erfahrungswerte auch bei Personalwechseln
erhalten bleiben, so dass die Qualität der Dienstleistungen und die
Effizienz der Arbeitsprozesse langfristig gesichert sind.

\section{Literatur}\label{literatur}

Pfadenhauer, Michaela (2002): Auf gleicher Augenhöhe reden - Das
Experteninterview - ein Gespräch zwischen Experte und Quasi-Experte, S.
113. In: Bogner, Alexander (Hrsg.): Das Experteninterview - Theorie,
Methode, Anwendung, Wiesbaden: Springer Fachmedien, 2002.

%autor
\begin{center}\rule{0.5\linewidth}{0.5pt}\end{center}

\textbf{Charlotte Lenz} hat
Bibliotheksmanagement an der Fachhochschule Potsdam sowie Bibliotheks-
und Informationswissenschaften an der HU Berlin studiert. Nach einer
langjährigen Anstellung als wissenschaftliche Bibliothekarin in der
Bibliothek des Deutschen Historischen Museums betreut sie nun die
Bibliothek der Gedenkstätte Sachsenhausen.

\textbf{Dr.~Matthias Miller} hat Historische Hilfswissenschaften,
Musikwissenschaften und Romanistik an der Universität Tübingen studiert.
Im Anschluss absolvierte er ein Referendariat an der Württembergischen
Landesbibliothek Stuttgart und der Bibliotheksschule Frankfurt/Main.
Nach \linebreak neun Jahren als wissenschaftlicher Mitarbeiter in der
Handschriftenabteilung an der Universitätsbibliothek Heidelberg übernahm
er 2008 die Leitung der Bibliothek am Deutschen Historischen Museum.

\end{document}